\chapter{КВАНТОВЫЕ ТОЧКИ И МОДЕЛИ НАНОСТРУКТУР В МЕТОДЕ СИЛЬНОЙ СВЯЗИ}

\section{Квантовое ограничение в низкоразмерных наноструктурах}

При переходе от массивного кристалла к наноструктуре электронные состояния начинают существенно зависеть от конечного размера образца. Если характерный размер области движения становится сравнимым с длиной волны электронного состояния, непрерывное описание в терминах объёмных зон уже не даёт полного представления о низкоэнергетической части спектра. Появляются уровни, положение которых определяется не только материалом и симметрией решётки, но и геометрией ограничивающей области~\cite{ashcroft1976solidstate,kittel2005solidstate}.

Такое изменение спектра связано с квантовым ограничением. В простейших моделях энергия низших состояний растёт при уменьшении характерного размера, причём ведущий масштаб часто имеет вид, близкий к обратному квадрату размера. Для реальных наноструктур эта зависимость может усложняться из-за атомной структуры границы, формы области, анизотропии и конкретных параметров материала. Тем не менее сама идея размерно-зависимого низкоэнергетического спектра остаётся базовой для физики низкоразмерных систем~\cite{reimann2002quantumdots,yoffe2001quantumdots}.

В задачах наноматериалов и нанотехнологий важна не только абсолютная величина основного уровня, но и расстояния между соседними уровнями. Спектральные зазоры определяют, насколько чувствительной система оказывается к изменению размера, формы и граничных условий. Поэтому анализ низших уровней конечной наноструктуры является естественной физической задачей даже в тех случаях, когда модель остаётся идеализированной.

Настоящая работа рассматривает именно такую идеализированную постановку. Цель не состоит в моделировании конкретного материала с параметрами из эксперимента или из первопринципного расчёта. Исследуется модельная квантовая точка на квадратной решётке, где можно контролируемо менять размер и форму границы и затем проверять, насколько простые физически мотивированные зависимости описывают рассчитанный спектр. Такая постановка удобна как промежуточный уровень между аналитическими учебными моделями и более тяжёлыми материал-специфическими расчётами.

\section{Квантовые точки как модельные наноструктуры}

Квантовую точку можно рассматривать как конечную область, в которой движение электронного состояния ограничено во всех пространственных направлениях. В результате низкоэнергетический спектр становится дискретным, а положение уровней зависит от формы и размера области~\cite{jacak1998quantumdots,kouwenhoven2001fewelectron}. В рамках данной работы квантовая точка понимается как модельная конечная наноструктура, а не как конкретный экспериментальный объект.

Физически такая модель полезна тем, что позволяет отделить геометрический фактор от других усложнений. Если фиксировать тип гамильтониана и менять только границу области, можно проследить, как низшие уровни реагируют на размер, отношение сторон и гладкость формы. Именно эта зависимость лежит в основе дальнейшего выбора признаков для суррогатной аппроксимации.

Идеальные прямоугольные и круговые области удобны для контроля, но они не исчерпывают возможные формы конечных наноструктур. Граница реального нанокристаллита или литографически заданной области редко является строго прямоугольной или строго круговой; размер и форма малых кристаллитов давно рассматриваются как факторы, влияющие на их электронные и оптические свойства~\cite{brus1984smallcrystallites,alivisatos1996quantumdots}. Поэтому полезна параметрическая семья форм, которая позволяет плавно переходить между более ромбоподобными, эллиптическими и более прямоугольноподобными границами, не меняя топологию области.

В данной работе такой семьёй служат суперэллипсы. Они дают простой способ управлять формой границы с помощью показателя $n$ и отношением сторон $r_{\mathrm{AR}}$. При этом в основной постановке используются фиксированные дискретные классы $n$, а не непрерывная модель произвольной формы. Такое ограничение делает задачу проверяемой и не требует утверждений об универсальном обобщении на все возможные границы.

С точки зрения физики наноструктур суперэллиптическая область может рассматриваться как модельная квантовая точка или конечный нанокристаллит с настраиваемой формой. В такой системе размер задаёт главный масштаб квантового ограничения, а форма определяет дополнительные расщепления и изменение расстояний между уровнями. Эти эффекты особенно заметны в низкоэнергетической части спектра, которая и является предметом дальнейшего анализа.

\section{Метод сильной связи и дискретные квантовые бильярды}

Метод сильной связи является одним из стандартных способов описания электронных состояний в кристаллической решётке~\cite{slater1954lcao,datta2005quantumtransport}. В этой постановке состояние задаётся амплитудами на узлах решётки, а движение между узлами описывается интегралами перескока. Такая модель сохраняет дискретность кристаллической структуры и поэтому отличается от чисто континуальных моделей эффективной массы, где решётка входит только через параметры непрерывного гамильтониана.

Для целей данной работы эта дискретность принципиальна. Квадратная решётка используется как минимальная контролируемая модель кристаллической дискретности. Она не означает, что исследуется конкретный квадратный материал, но позволяет явно видеть, как конечная граница, число узлов и симметрии решётки влияют на уровни энергии. Поэтому остаточные решёточные эффекты не считаются численной ошибкой сами по себе: они являются частью выбранной физической модели.

Конечные области решётки с открытыми границами естественно связаны с понятием дискретного квантового бильярда. В континуальном квантовом бильярде частица ограничена областью с жёсткой стенкой; в дискретном варианте область задаётся конечным множеством узлов решётки, а перескоки за границу отсутствуют. Спектр такой системы зависит одновременно от общей формы области и от того, как её граница ложится на узлы решётки. Квантовые бильярды используются как модельные системы для анализа спектральной статистики и влияния формы границы~\cite{berry1981sinai,stockmann1990quantumchaos,cuevas1996chaoticbilliards}. Бильярды в модели сильной связи обсуждаются как отдельный класс конечных решёточных систем~\cite{ulcakar2022tightbindingbilliards}; близкие идеи возникают и в контексте квантовых бильярдов на оптических решётках~\cite{montangero2008quantumbilliards}.

В отличие от континуальной задачи, дискретная область имеет дополнительные геометрические характеристики. Например, фактическое число узлов может немного отличаться от аналитической площади гладкой фигуры, а квадратная решётка допускает разбиение на две подрешётки. Такие особенности могут влиять на отдельные уровни или на остатки аппроксимации, особенно для малых размеров и более резких границ. Поэтому в последующих главах проверяются не только основные размерные зависимости, но и диагностические величины, связанные с дискретизацией границы.

Численная реализация таких конечных систем выполняется с помощью Kwant~\cite{groth2014kwant}. Однако физический объект в работе задаётся не программным пакетом, а гамильтонианом метода сильной связи и выбранной конечной геометрией. Kwant используется как инструмент построения конечной решёточной системы и прямого расчёта её спектра. Подробная запись гамильтониана и расчётной процедуры приведена в главе 2.

\section{Суррогатные модели в вычислительной физике наноструктур}

Прямой расчёт спектра конечной квантовой системы является основной физической операцией в данной работе. В более сложных задачах, особенно при переборе многих геометрий или при переходе к материал-специфическим моделям, многократные прямые расчёты могут становиться трудоёмкими. Поэтому в вычислительной физике и инженерном моделировании используются суррогатные аппроксимации: они приближают уже рассчитанную зависимость и позволяют быстро исследовать область параметров~\cite{queipo2005surrogate,forrester2008surrogate,hou2024surrogateprotocol}.

В настоящей работе суррогатная модель не рассматривается как замена квантового расчёта. Она не создаёт новую физику и не исправляет исходный гамильтониан. Её задача уже: по набору прямых расчётов Kwant проверить, можно ли компактно описать низкоэнергетический спектр через физически мотивированные геометрические величины. Поэтому качество суррогатной аппроксимации всегда оценивается относительно прямого спектрального расчёта.

Такой подход особенно естественен, когда признаки выбираются не произвольно, а из ожидаемых физических масштабов. Использование научного знания в моделях машинного обучения рассматривается как отдельная методологическая линия, но в настоящей работе оно применяется в узком физическом смысле: признаки задаются низкоэнергетическим масштабированием, а не подбираются слепым поиском~\cite{willard2022piml}. Если низшая энергия, отсчитанная от дна зоны, находится в квадратичном низкоэнергетическом режиме, то размерная зависимость порядка $1/a^2$ должна проявляться уже в простой модели. Тогда интерпретируемая линейная аппроксимация может оказаться не менее ценной, чем более сложная нелинейная модель, потому что она непосредственно связана с физическим механизмом квантового ограничения.

В дальнейшем в работе сравниваются гребневая регрессия (Ridge) и многослойный перцептрон (МП) малой ёмкости. Это сравнение не предназначено для демонстрации преимуществ нейросетевых методов как таковых. Оно отвечает более узкому физическому вопросу: достаточно ли признаков, основанных на низкоэнергетическом масштабировании, чтобы аппроксимировать $E_0$ и $dE_1$ в проверенном диапазоне размеров и форм, или добавление малой нелинейности даёт устойчивый выигрыш. Примеры применения машинного обучения к задачам квантовых точек существуют и в более материал-ориентированных постановках~\cite{zielinski2025nanowireml}; здесь же нейросетевая модель используется только как контрольная аппроксимация.

Материал-специфическое развитие такой методики возможно, например, через калибровку эффективных параметров метода сильной связи или через построение отдельной эталонной выборки с использованием теории функционала плотности (ТФП) и OpenMX~\cite{hohenberg1964inhomogeneous,kohn1965selfconsistent,ozaki2003openmx}. В данной дипломной работе этот этап не выполняется. Рассматриваемая модель остаётся контролируемой решёточной постановкой, предназначенной для проверки связи между геометрией модельной наноструктуры и её низкоэнергетическим спектром.

\section{Выводы по главе 1}

В этой главе сформулирована физическая мотивация работы. Основные выводы состоят в следующем.

\begin{enumerate}
  \item В низкоразмерных наноструктурах конечный размер и граничные условия существенно влияют на низкоэнергетический спектр из-за квантового ограничения.
  \item Модельные квантовые точки и конечные нанокристаллиты являются удобными объектами для изучения зависимости спектральных уровней от размера и формы границы.
  \item Метод сильной связи позволяет сохранить дискретность кристаллической решётки и тем самым исследовать не только континуальное размерное масштабирование, но и решёточные эффекты границы.
  \item Конечные области решётки с открытыми границами связаны с дискретными квантовыми бильярдами, где спектр определяется как формой области, так и её дискретным представлением.
  \item Суррогатная аппроксимация в данной работе используется как вспомогательный и проверяемый инструмент, а не как замена прямого расчёта спектра.
  \item Эти соображения приводят к модели и численной методике главы 2: квадратная решётка метода сильной связи, конечные суперэллиптические области и прямой расчёт низших уровней в Kwant.
\end{enumerate}

\clearpage
